\documentclass[11pt]{article}
\usepackage{amsmath}

\setlength{\parindent}{0pt}
\setlength{\parskip}{0.7em}
\setlength{\textwidth}{6.5in}
\setlength{\oddsidemargin}{0in}
\setlength{\evensidemargin}{0in}
\setlength{\topmargin}{-0.4in}
\setlength{\textheight}{9.2in}

\begin{document}

\begin{center}
{\Large SYSC 5108 Exam Study: Assignment 2, Question 2}\\
\vspace{0.3em}
{\large MLE for the Laplace rate parameter}\\
\vspace{0.3em}
Source question: \texttt{SYSC 5108 W26 Assignment 2.pdf}, Question 2
\end{center}

\section*{Question}

Given $n$ independent sample points $x_i$, $i=1,\ldots,n$, from a Laplace distribution with PDF
\[
f(x)=\frac{\lambda}{2}e^{-\lambda |x|},
\qquad \lambda>0,
\]
provide the maximum likelihood estimator (MLE) for the parameter $\lambda$.

\section*{Step 1: Identify the model and parameter}

The PDF is a zero-centered Laplace distribution:
\[
f(x;\lambda)=\frac{\lambda}{2}e^{-\lambda |x|}.
\]
The unknown parameter is $\lambda$, which controls how sharply the density is concentrated around zero. Larger $\lambda$ gives a sharper peak and lighter tails.

Goodfellow, Bengio, and Courville define the Laplace distribution in Chapter 3 as
\[
\operatorname{Laplace}(x;\mu,\gamma)
=
\frac{1}{2\gamma}\exp\left(-\frac{|x-\mu|}{\gamma}\right).
\]
The assignment's distribution is the special case
\[
\mu=0,\qquad \gamma=\frac{1}{\lambda}.
\]
So estimating $\lambda$ is equivalent to estimating the inverse scale parameter.

\section*{Step 2: Write the likelihood}

The samples are independent, so the joint likelihood is the product of the individual densities:
\[
L(\lambda)
=
\prod_{i=1}^{n} f(x_i;\lambda).
\]
Substitute the PDF:
\[
L(\lambda)
=
\prod_{i=1}^{n}\left(\frac{\lambda}{2}e^{-\lambda |x_i|}\right).
\]
Collect the repeated factors:
\[
L(\lambda)
=
\left(\frac{\lambda}{2}\right)^n
\exp\left(-\lambda\sum_{i=1}^{n}|x_i|\right).
\]

\section*{Step 3: Take the log-likelihood}

MLE usually maximizes the log-likelihood because the logarithm turns products into sums and does not change the maximizer. This is the same principle used in the Chapter 3 slides on MLE and in Goodfellow et al., Chapter 5, Section 5.5.

The log-likelihood is
\[
\ell(\lambda)
=
\log L(\lambda)
=
n\log\lambda-n\log 2-\lambda\sum_{i=1}^{n}|x_i|.
\]

\section*{Step 4: Differentiate and solve the score equation}

Differentiate with respect to $\lambda$:
\[
\frac{d\ell}{d\lambda}
=
\frac{n}{\lambda}
-
\sum_{i=1}^{n}|x_i|.
\]
Set the derivative equal to zero:
\[
\frac{n}{\lambda}
-
\sum_{i=1}^{n}|x_i|
=0.
\]
Rearrange:
\[
\frac{n}{\lambda}
=
\sum_{i=1}^{n}|x_i|.
\]
Therefore
\[
\boxed{
\hat{\lambda}_{\mathrm{MLE}}
=
\frac{n}{\sum_{i=1}^{n}|x_i|}
}.
\]

Equivalently,
\[
\boxed{
\hat{\lambda}_{\mathrm{MLE}}
=
\frac{1}{\frac{1}{n}\sum_{i=1}^{n}|x_i|}
}.
\]
So the MLE is the inverse of the sample mean absolute value.

\section*{Step 5: Check that it is a maximum}

The second derivative is
\[
\frac{d^2\ell}{d\lambda^2}
=
-\frac{n}{\lambda^2}.
\]
For $\lambda>0$,
\[
-\frac{n}{\lambda^2}<0.
\]
So the stationary point is a maximum.

One edge case: if every $x_i=0$, then $\sum_i |x_i|=0$, and the likelihood increases without bound as $\lambda\rightarrow\infty$. In that degenerate case there is no finite MLE. For ordinary data with at least one nonzero sample,
\[
\sum_{i=1}^{n}|x_i|>0,
\]
and the finite MLE above is valid.

\section*{Step 6: Exam Interpretation}

The negative log-likelihood is
\[
-\ell(\lambda)
=
-n\log\lambda+n\log 2+\lambda\sum_{i=1}^{n}|x_i|.
\]
This shows why the Laplace distribution is tied to absolute error: the data-dependent term is proportional to
\[
\sum_{i=1}^{n}|x_i|.
\]
This contrasts with the Gaussian example from the Chapter 3 slides and Goodfellow et al., Chapter 5, where Gaussian maximum likelihood leads to squared-error terms.

A concise exam answer could be:
\begin{quote}
Since the samples are independent,
\[
L(\lambda)
=
\prod_{i=1}^{n}\frac{\lambda}{2}e^{-\lambda |x_i|}
=
\left(\frac{\lambda}{2}\right)^n
e^{-\lambda\sum_i |x_i|}.
\]
Thus
\[
\ell(\lambda)=n\log\lambda-n\log 2-\lambda\sum_i |x_i|.
\]
Setting the score to zero gives
\[
\frac{d\ell}{d\lambda}
=
\frac{n}{\lambda}-\sum_i |x_i|
=0,
\]
so
\[
\hat{\lambda}_{\mathrm{MLE}}
=
\frac{n}{\sum_i |x_i|}.
\]
The second derivative is $-n/\lambda^2<0$, confirming a maximum.
\end{quote}

\section*{Cross References Used}

\begin{itemize}
\item Assignment: \texttt{SYSC 5108 W26 Assignment 2.pdf}, Question 2.
\item Local submitted Assignment 2 PDF checked for comparison, Question 2 response.
\item Slides: Chapter 3 slide deck, slide 26 on the MLE principle and maximizing the log-likelihood of independent samples.
\item Slides: Chapter 3 slide deck, slide 27 on the Gaussian MLE example, useful as a contrast to the Laplace absolute-error form.
\item Slides: Chapter 3 slide deck, slide 28 on MLE properties.
\item Textbook: Goodfellow, Bengio, and Courville, \textit{Deep Learning}, Chapter 3, Section 3.9.4 on the Laplace distribution.
\item Textbook: Goodfellow, Bengio, and Courville, \textit{Deep Learning}, Chapter 5, Section 5.5 on maximum likelihood estimation and log-likelihood.
\end{itemize}

\end{document}
